\item \textbf{Định nghĩa sự kiện lập lịch (\textit{Scheduling Event}):} Nhóm quy ước một sự kiện lập lịch của một tiến trình là một lần tiến trình đó được \texttt{scheduler()} chọn từ trạng thái \texttt{RUNNABLE} và thực hiện chuyển ngữ cảnh (\texttt{swtch()}) để trao quyền thực thi trên CPU. Việc chỉ tăng \texttt{p->nsched++} tại một vị trí duy nhất trong \texttt{scheduler()} ngay trước \texttt{swtch()} giúp logic tập trung (single point of truth), tránh can thiệp rải rác vào nhiều hàm nhường CPU khác nhau (\texttt{yield}, \texttt{sleep}, \texttt{exit}) và loại trừ hoàn toàn nguy cơ đếm trùng lặp.

\item \textbf{Chiến lược sao chép từng phần tử thay vì toàn bộ cấu trúc (\textit{Streaming Copyout}):} Toàn bộ \texttt{struct sysinfo} có kích thước $\approx 38.25\text{ KB}$, trong khi kernel stack chỉ có đúng $1$ trang $4\text{ KB}$ (\texttt{PGSIZE = 4096}). Để triệt tiêu nguy cơ tràn kernel stack mà không phải cấp phát động phức tạp hay làm cạn kiệt bộ nhớ trang qua \texttt{kalloc()}, hàm \texttt{proccollect()} chỉ khai báo biến cục bộ cho từng phần tử \texttt{struct sysinfo\_proc} ($160$ bytes), đóng gói dữ liệu rồi lập tức ghi vào không gian user bằng \texttt{copyout()} thông qua tính offset chính xác bằng \texttt{\_\_builtin\_offsetof}.

\item \textbf{Thứ tự khóa (\textit{Lock Ordering}) chống Deadlock:} Để lấy \texttt{ppid}, ta cần truy xuất \texttt{p->parent->pid} vốn được bảo vệ bởi \texttt{wait\_lock}. Trong xv6, ở các hàm như \texttt{exit()} và \texttt{wait()}, \texttt{wait\_lock} luôn được lấy trước \texttt{p->lock}. Do đó, \texttt{proccollect()} tuân thủ nghiêm ngặt: lấy \texttt{wait\_lock} trước, sau đó mới lấy \texttt{p->lock}, và giải phóng theo thứ tự ngược lại để ngăn chặn tình trạng khóa chết vòng tròn (circular wait deadlock).

\item \textbf{Giả định về việc đọc bảng \texttt{ofile[]} trong \texttt{proccollect}:} Trong xv6, mảng \texttt{ofile[NOFILE]} là thuộc tính private của tiến trình. Khi \texttt{proccollect()} duyệt qua \texttt{ofile[]} của tiến trình khác, thao tác được bảo vệ dưới khóa \texttt{p->lock}. Thiết kế chủ động không giữ \texttt{ftable.lock} trong khi đang giữ \texttt{p->lock} để tránh tranh chấp thứ tự khóa gây deadlock trên hệ thống đa lõi; kết quả đọc là một bản chụp thời điểm (point-in-time snapshot) hoàn toàn chấp nhận được cho mục đích chẩn đoán hệ thống.
